¿Cuánto podemos fiarnos de una mediana? Esta gráfica compara, para cada medida, la ronda más lenta con la más rápida. Uno significa que todas las rondas coinciden.

Sin contención, todo está en uno: el cronómetro marca siempre lo mismo. Con 24 hilos, la misma medida cambia típicamente entre 3 y 6 veces de una ronda a otra, y en el peor caso, 37 veces.

¿Por qué con 15 rondas sale peor que con 5? Porque con más rondas hay más ocasiones de pillar una muy buena y una muy mala, como al tirar un dado más veces. Más rondas no hacen cada ronda más estable; hacen la mediana menos dependiente de la mala suerte.

Consecuencia: con contención, una diferencia del 20 por ciento entre dos algoritmos no significa nada. Solo valen las enormes, como la de Libstdc++.
